% HD-G14. 21 November 2026. GPT-6 Astra Pro, inherited effort.
\section{An actual nonsimple fourfold seed and the kernel of its square}

\paragraph{HD-G14.}
Let $X$ be a complex abelian fourfold, let $p\in X$, and put
\[
 E=\mathcal O_X\oplus\mathcal O_p,\qquad
 F=E\boxtimes E\quad\hbox{on }X\times X.
\]
Then $E$ is semiregular and has no negative self-Ext.
Its endomorphism space has dimension three.
Its external square has the following exact degree-two trace kernel:
\[
 \ker\sigma_F^2=
 \bigl(\operatorname{Ext}^2(E,E)\otimes\ker\sigma_E^0\bigr)
 \oplus
 \bigl(\ker\sigma_E^0\otimes\operatorname{Ext}^2(E,E)\bigr),
 \qquad \dim\ker\sigma_F^2=24.                         \tag{G14.1}
\]
Both $E$ and $F$ have trivial stabilizer under translations and
tensor products by degree-zero line bundles. In particular,
no finite point subgroup acting on this $F$ makes it weakly
equivariantly semiregular.

\paragraph{The Ext groups and every trace map needed here.}
The local Koszul resolution at $p$ gives
\[
 \operatorname{Ext}^j(\mathcal O_p,\mathcal O_p)
       =\bigwedge^jT_pX.
\]
The other three blocks are
\[
 \operatorname{Ext}^j(\mathcal O_X,\mathcal O_X)
       =H^j(X,\mathcal O_X),\qquad
 R\operatorname{Hom}(\mathcal O_X,\mathcal O_p)=\mathbb C[0],
 \qquad
 R\operatorname{Hom}(\mathcal O_p,\mathcal O_X)=\mathbb C[-4].
                                                               \tag{G14.2}
\]
The last equality follows either by dualizing the same Koszul
resolution or by Serre duality, with a chosen holomorphic volume
form. Thus the dimensions of $\operatorname{Ext}^j(E,E)$ for
$j=0,1,2,3,4$ are respectively
\[
                 3,\ 8,\ 12,\ 8,\ 3,                       \tag{G14.3}
\]
and all other degrees vanish. In particular $\chi(E,E)=2$.

The Atiyah class of a direct sum is block diagonal. Composition
with any off-diagonal endomorphism has zero trace after every
Atiyah power. Consequently the nonzero map
$\xi:\mathcal O_X\to\mathcal O_p$, viewed as an off-diagonal
endomorphism of $E$, satisfies $\sigma_E^0(\xi)=0$.
The two diagonal identity endomorphisms have trace
$1$ and $[p]$ respectively; these lie in different Hodge
summands and are linearly independent. Therefore
\[
                    \ker\sigma_E^0=\mathbb C\xi.          \tag{G14.4}
\]

For $\mathcal O_X$, the positive-degree trace is simply the
identity on $H^j(X,\mathcal O_X)$, since its Atiyah class
vanishes. For $\mathcal O_p$, global tangent vectors identify
with $T_pX$ by translation. The local Koszul calculation of
the Hochschild action identifies
\[
 H^0(X,\bigwedge^jT_X)\longrightarrow
             \operatorname{Ext}^j(\mathcal O_p,\mathcal O_p)
\]
with the identity on $\bigwedge^jT_pX$.
The trace-evaluation formula of HD-G13 and HD-N11 therefore
identifies its composite trace with contraction against $[p]$.
Writing a nonzero top class as the product of four holomorphic
and four antiholomorphic one-forms proves directly that this
contraction is an isomorphism onto
$H^4(X,\Omega_X^{4-j})$ for $0\le j\le4$.
For $j=1,2$, this Hodge summand is disjoint from
$H^j(X,\mathcal O_X)$. There are no off-diagonal Ext groups
in these degrees by(G14.2). Hence
\[
       \sigma_E^1\ \hbox{and}\ \sigma_E^2
                     \ \hbox{are injective}.              \tag{G14.5}
\]
This proves the claimed semiregularity of $E$. It makes no
assertion of trace injectivity in degree zero or in every
degree.

\paragraph{The exact external-product calculation.}
The Ext K\"unneth isomorphism and the absence of negative Ext
identify $\operatorname{Ext}^2(F,F)$ with the direct sum
\[
 \operatorname{Ext}^2(E,E)\otimes\operatorname{End}(E)
 \ \oplus\
 \operatorname{Ext}^1(E,E)\otimes\operatorname{Ext}^1(E,E)
 \ \oplus\
 \operatorname{End}(E)\otimes\operatorname{Ext}^2(E,E).
                                                               \tag{G14.6}
\]
For external products,
$\operatorname{At}(E\boxtimes E)
 =\operatorname{At}(E)\boxtimes1+1\boxtimes\operatorname{At}(E)$.
The divided exponential and the graded tensor trace then
give the tensor product of the corresponding trace maps
on each summand, with the usual fixed Koszul sign.
The three images lie respectively in the distinct Hochschild
K\"unneth summands
\[
 HH_{-2}(X)\otimes HH_0(X),\qquad
 HH_{-1}(X)\otimes HH_{-1}(X),\qquad
 HH_0(X)\otimes HH_{-2}(X).                                \tag{G14.7}
\]
They cannot cancel across these summands.
Equations(G14.4)--(G14.5) now give exactly(G14.1).
In particular $\operatorname{ext}^2(F,F)=136$ and
$\operatorname{rank}\sigma_F^2=112$.

\paragraph{The actual point stabilizers and Fourier transport.}
The rank-one sheaf $E$ has determinant $\mathcal O_X$.
For completeness, the point sheaf has trivial determinant
because its support has codimension at least two: its
determinant line is canonically trivial off that support,
and a divisor cannot be supported in codimension at least two.
Thus an isomorphism $t_a^*E\otimes P\simeq E$ forces
$P\simeq\mathcal O_X$ by taking determinants.
The torsion subsheaf of $E$ is intrinsically $\mathcal O_p$,
so its support then forces $a=0$.

Likewise $F$ has rank one, determinant $\mathcal O_{X\times X}$,
and torsion support exactly
\[
             (X\times\{p\})\ \cup\ (\{p\}\times X).
                                                               \tag{G14.8}
\]
Taking determinants again forces the tensor line to be
trivial. A translation by $(a,b)$ sends the two irreducible
components to the same coordinate types with point
coordinates $p-b$ and $p-a$. Equality of the supports
forces $a=b=0$. This proves the full point-stabilizer
assertion without choosing a decomposition of an isomorphic
sheaf. A finite point subgroup admitting an equivariant
structure on $F$ is consequently trivial; its invariant
category and its trace map are the original ones.

Coordinate isomorphisms, tensor products by line bundles,
and partial Poincare Fourier transforms used in the seed
construction transport these point stabilizers by actual
invertible maps. If $\Phi_{a,P}=t_a^*(-)\otimes P$, then
conjugation by tensoring with $L$ replaces $P$ by
$P\otimes L\otimes(t_a^*L)^{-1}$. Pullback by an isomorphism
$u$ sends $(a,P)$ to $(u^{-1}a,u^*P)$.
For the positive Poincare kernel, the change of integration
variable $z=x+a$ and the biextension identity give the
point-coordinate map $(a,b)\mapsto(b,-a)$; the inverse
kernel gives its inverse. Each is an isomorphism on the
translation/tensor point group. The corresponding enhanced
equivalences identify Ext groups and induce isomorphisms on
Hochschild homology, commuting with the trace.
Therefore the nonzero kernel(G14.1), and the absence of a
nontrivial point stabilizer, persist through these
coherently conjugated Fourier constructions.
This assertion concerns these explicit operations; it does
not identify arbitrary new equivariant structures after
unrelated changes of the object.

\paragraph{Relation to the exact character obstruction.}
Here $\operatorname{ch}(E)=1+[p]$. This character cannot have
the form in HD-N11 for any ample $H$, algebraic rational $B$,
and $s^2=-d<0$. Indeed its rank gives
$\alpha+\bar\alpha=1$. Set
$u=s(\alpha-\bar\alpha)\in\mathbb Q$.
The vanishing of its degree-two cohomology component gives
$B=-uH$. Its degree-four component would then be
\[
                 \frac{s^2-u^2}{2}H^2
                =-\frac{d+u^2}{2}H^2\ne0,
                                                               \tag{G14.9}
\]
whereas that component of $1+[p]$ is zero.
The inequality follows from ampleness, which makes $H^2$
nonzero, and $d>0$.
Thus this is an actual nonsimple semiregular fourfold object,
with an explicit character comparison and an explicit
external-square obstruction. It supplies neither the
HD-N11 character nor a semiregular external-square source
for the remaining eightfold construction.
\hfill$\square$
